\section{תקשורת בין תוכנות: SDK, GUI, CLI, API}

תוכנות מודרניות נבנות בשכבות. הבנת השכבות חיונית לתכנון סוכנים שצריכים
לדבר עם מערכות קיימות --- בין אם דרך ממשק אנושי או ישירות דרך ממשק מכונה.

% Software communication stack
\begin{figure}[H]
\centering
\begin{tikzpicture}[
  layer/.style={draw=secondary, fill=white, minimum width=8.5cm,
    minimum height=0.85cm, align=center, font=\small}
]
  \node[layer, fill=accentamber!12] (api) at (0,3.2) {\eng{API} --- חשיפה חיצונית לתוכנה};
  \node[layer, fill=accentblue!10] (cli) at (0,2.2) {\eng{CLI} --- ממשק טקסטואלי בטרמינל};
  \node[layer, fill=accentgreen!10] (gui) at (0,1.2) {\eng{GUI} --- ממשק גרפי למשתמש אנושי};
  \node[layer, fill=primary!10] (sdk) at (0,0.2) {\eng{SDK} --- שכבת פיתוח מעל הפונקציות};
  \node[layer, fill=lightgray] (fn) at (0,-0.8) {פונקציות תוכנה --- ליבת הביצוע};
  \foreach \a/\b in {fn/sdk, sdk/gui, sdk/cli, sdk/api}{
    \draw[-{Stealth[length=2mm]}, thick, primary] (\a.north) -- (\b.south);
  }
\end{tikzpicture}
\caption{מחסנית תקשורת בין תוכנות: פונקציות, \eng{SDK}, \eng{GUI}/\eng{CLI} ו-\eng{API}}
\label{fig:software-stack}
\end{figure}


\subsection{שכבת הפונקציות}

בתחתית נמצאות פונקציות התוכנה: חישוב, חיפוש, כתיבה לדיסק. שכבה זו אינה
נגישה ישירות למשתמש הקצה --- היא ליבת הלוגיקה.

\subsection{SDK --- ערכת פיתוח}

\eng{SDK} (\eng{Software Development Kit}) עוטף את הפונקציות בממשק נוח
למפתחים. כל פעולה חיצונית בתוכנה אמורה לעבור דרך ה-\eng{SDK}. כשבונים
\eng{GUI} או \eng{CLI}, הם קוראים ל-\eng{SDK} ולא ישירות לפונקציות גולמיות.

\subsection{GUI --- ממשק גרפי}

\eng{GUI} (\eng{Graphical User Interface}) מיועד לבני אדם: תפריטים, חלונות,
גרירה בעכבר. דוגמאות: PowerPoint, Excel, דפדפן. סוכן שעובד דרך \eng{GUI}
מדמה משתמש אנושי --- לוחץ, מצלם מסך, מקליד --- וזה יקר בטוקנים ובזמן.

\subsection{CLI --- ממשק שורת פקודה}

\eng{CLI} (\eng{Command Line Interface}) הוא תפריט טקסטואלי בטרמינל:
לחיצה על \eng{A} מבצעת פעולה אחת, \eng{B} --- אחרת. פשוט, מהיר וידידותי
לסוכנים. ההרצאה מדגישה מעבר מעבודה ב-\eng{GUI} לעבודה בטרמינל ובטקסט.

\subsection{API --- חשיפה לעולם החיצון}

\eng{API} (\eng{Application Programming Interface}) מאפשר לתוכנה חיצונית
לבצע פעולות מורשות. כשתוכנה א' רוצה לדבר עם תוכנה ב', הן מתקשרות דרך
\eng{API} --- בתנאי שיודעים לתכנת את החיבור.

\begin{lstlisting}[language=Java,caption={דוגמת בקשת API (JSON)}]
POST /api/v1/documents/compile HTTP/1.1
Content-Type: application/json

{
  "source": "main.tex",
  "engine": "lualatex",
  "runs": 3
}
\end{lstlisting}

\agenttable{
\caption{השוואת דרכי גישה לתוכנה}
\label{tab:access-methods}
\begin{tabular}{@{}p{2cm}p{2.5cm}p{3.5cm}p{3.5cm}@{}}
\toprule
\textbf{שיטה} & \textbf{קהל יעד} & \textbf{יתרון} & \textbf{חיסרון לסוכן} \\
\midrule
\eng{GUI} & אדם & אינטואיטיבי & טוקנים, איטיות \\
\eng{CLI} & מפתח/סוכן & מהיר, טקסט & דורש למידה \\
\eng{API} & מערכות & אוטומציה מלאה & דורש תיעוד \\
\eng{SDK} & מפתח פנימי & שליטה מלאה & לא חשוף תמיד \\
\bottomrule
\end{tabular}
}

\subsection{Skill כעטיפה ל-API}

כפי שנלמד בהרצאה, ניתן לעטוף \eng{API} ב-\eng{Skill} ולאפשר ``\eng{AI via API}''
--- שפה חופשית שמתורגמת לקריאות מובנות. כך תוכנה א' ותוכנה ב' יכולות
לתקשר דרך סוכנים בלי שהמפתח יזכור כל פרמטר.

\subsection{כיוון עתידי}

החזון מההרצאה: מערכות הפעלה מבוססות \eng{AI} שבהן המשתמש מדבר בשפה
חופשית והמערכת בוחרת כלים (\eng{Excel}, Python, דפדפן) בשקיפות. למימוש
חזון זה נדרש שרוב התוכנות יחשפו יכולות בטקסט/\eng{API}, לא רק ב-\eng{GUI}.
